\documentclass{article}
\usepackage[margin=2cm]{geometry}
\usepackage{amsmath}
\usepackage{graphicx}
\usepackage{booktabs}
\usepackage[utf8]{inputenc}
\usepackage{ragged2e}
\setlength{\emergencystretch}{3em}
\begin{document}
\sloppy

(2) A right-of-way is reserved for
ditches and canals constructed byauthority of the United States under the
Act of August 30, 1890 (43 U.S.C. 945).
(3) The parcels are subject to valid
existing rights.
and food ol a.oareservations for roads, public utilities,existing and proposed, in accordancewith the local governing entities(4) The parcels are subect to
transportation plans.


The following numbered terms and
conditions will appear on theconveyance documents for the sale
parcels:


(5) An appropriate indemnification
Claims arising out of the patentee's use,clause protecting the United States from
palented lands.occupancy, or occupations on the


parcels are subject to the roquirementsof Section 120(h) of the CERCLA, asTo the extent required by law. the
amended. Accordingly, notice is hereby
given that the lands have been
subject properties.have been stored for one yoar or morenor that any hazardous substances havebeen disposed of or released on the


No warranty of any kind, express or
circmstanc or condition. Tha contingency basis.the land may be developed, its physicalcondition utre uses, or any otherconveyance of the arcels will not be on


The regulation provides that allother
use, absent statutory er other expressauthority, requires a sales contract or
permit. The BLM refers interestedpartie to the explanation of this
Register in 2001, available at htp:/
not include large-scale use of mineralmaterials, oven within the boundaries the surface estate [66 FR 58894].
Further explanation is contained in the
BLM Instruction Memorandum No.
the BLM's webait at hps:/www.blm.gov/policy/im-2014-085.2014085 (April 23, 2014), available on


(1) All mineral daposits in the lands
for, mine, and remove such deposits
Secretary ae reserved to the UnitedStates, together with all necessary
access and exit rights.


5055647689 or BIA Navajo RegionBLM Farmington Field Office ProjectManager Sarh Scott, scotrabm.govOffice Regional Archeologist/ProjectManager Robert Begay, robert.begay10bia.gov. 5058638515. Individuals indeafolind, hard of hearing or have aFOR FURTHERINFORMATION CONTACT:the United States who are deaf.
telecommunications melay servioe for
should use the relay services offered
contact in the United Staies.within their country to makeintenational calls to the point-of-


Termination of Preparation of the
Environmental Impact Statement forthe Farmington Mancos-GallupResource Management Plan
Amendment, New Mexico


SUMmany: The Bureau of Land
Management (BLM) and the Bureau of
ptio l actIndian Affi (BIA) ae teminating thestatement (EIS) Jor the Farmington
Mancos-Gallup Resource ManagementPlan (RMP) Amendment.
DATES: The EIS development process fotAmendment is terminated immediately.the Farmington Mancos-Gallup RMP


Dr yD 
FRD 2112 711-24, 845 ]
BILLING C0DE 431-21-P


Bureau of Land Management
[BLM\_NM\_FRN\_MO4500178179]


\subsection*{D. EPARTMENT OF THE INTERIOR}


SUPPLEMENTARY INFORMATION: Pursuant
to the National Envizonmental Policy
Act of 1969, as implemented by theCouncil on Environmental Quality
and associated EIS on February 25, 2014
(79 FR 10548), On October 21, 2016, the
BLM and the BIA published anamended Notice of Intent in the FederalRegister announcing the addition of the
1R2BIA as a joint/co-lead agency for the EISw to analy the impacts of additionaloil and gas development within the SanJuan Basin in northwestern New


AGENCY: Buresu of Land Management;
Burean of Indian Affairs, Interior.
Acnore Notice of termination.


\subsection*{Docket 1BSEE20240003: 00000
245E1700D2 ET1SF0000.EAQ000; OMB
Control Number 10140023]}


\subsection*{D. EPARTMENT OF THE INTERIOR}


Melanie G. Barnes,
Deberah s. Shirey.BMNwecic SDi
Actig o Rn Dr.
FRDoc 22435278 Fed7126;845 m)
BILLING C0DE 431-23-P


AGENCY: Bureau of Safety and
Environmental Enforcement, Interior.


(Authority: 40 CFR 1506.6, 40 CFR 1506.16)


Melanie G. Barnes,


\subsection*{Bureau of Safety and Environmental
Enforcement}


Mexico, as well as decisions related to
lands and realty.BLM-managed lands
with wilderness characteristics, and
also to evaluate alternatives and issuesvegetation management. The EIS was
related to the BIAs authority over
mineral leasing and associated activity
decisions on Navajo Tribal Trust Lands
and Navajo Indian allotments. The
Notice of Availability for the Draft EIS
published in the Federal Register on
February 28, 2020 (85 FR 12012).Thebureaus distributed the Draft EIS to
various Federal, State, and localagencies, elected officials, special
interest groups, interested individuals,
and the media.Due to the COVID-19
pandemic and restrictions placed on inperson meetings. virtual public hearings
were held on May 14, 15, 16, and 18
2020, as well as on August 26, 27, 28,and 292020.Since the initial
publication of the Notices of Intent in2014 and 2016, and the publication of
the draft RMP Amendment and EIS in
2020, there have been many changes
relevant to the plan amendment and
associated EIS, such as a change in thedevelopment trends in the San Juan
Basin; the withdrawal of 336,404 acres
from mineral entry around the Chaco
Culture National Historical Park; the
studies for the region; the establishmentof the Honoring Chaco Initiative; and anpreparation of BIA-funded ethnographic
increase in outdoor recreation in the
region.Given these changes and the
extent of revisions necessary to address
these changes in the current EIS
process, the agencies determined it is
impractical to continue the plan
amendment effort as currently
structured. Therefore, the BLM and BIA
hereby terminate preparation of the EIS
for the RMP Amendment.


Agency Information Collection
Activities; Pollution Prevention and
Control

\end{document}
